\documentclass[11pt]{article}
\usepackage[margin=1in]{geometry}
\usepackage{amsmath,amssymb,amsthm}
\usepackage{hyperref}
\usepackage{enumitem}

\newtheorem{theorem}{Theorem}
\newtheorem{lemma}[theorem]{Lemma}
\newtheorem{corollary}[theorem]{Corollary}
\theoremstyle{definition}
\newtheorem{definition}[theorem]{Definition}
\theoremstyle{remark}
\newtheorem{remark}[theorem]{Remark}

\usepackage{xurl}
\let\oldus\_
\renewcommand{\_}{\oldus\allowbreak}

\title{A disproof of Conjecture 00000001213\\[2pt]
\large In Sylver coinage the prime $5$ is a winning first move}
\date{}

\begin{document}
\sloppy
\maketitle

\section{The conjecture}

Conjecture 00000001213 reads:

\begin{quote}
\emph{Definition:} Sylver Coinage: two players alternately name positive integers that may not be
nonnegative integer combinations of previously named ones; naming 1 loses. \emph{Conjecture:} The
first player wins, and every winning first move is composite; no prime first move is winning ---
after a prime opening the second player has a winning response.
\end{quote}

The conjecture makes three claims:
\begin{itemize}[nosep]
\item (C1) the first player wins;
\item (C2) every winning first move is composite;
\item (C3) no prime first move is winning, i.e.\ after a prime opening the second player has a winning
response.
\end{itemize}
We prove that (C1) is true and that (C2) and (C3) are false. Indeed, every prime $p \ge 5$ is a winning
first move. This is Hutchings' theorem; we give a complete proof and formalize it in Lean 4.

\section{Definitions}

A \emph{position} is the finite set $S$ of numbers named so far. A number $n$ is \emph{legal} after $S$
if $n > 0$ and $n \notin \langle S \rangle$, where $\langle S\rangle$ is the additive submonoid of
$\mathbb{N}$ generated by $S$, i.e.\ the set of nonnegative integer combinations of $S$. Game values are
defined inductively, as in the accepted solution of Conjecture 00000008869:
\begin{itemize}[nosep]
\item the player to move from $S$ \emph{wins} if some legal $n \ne 1$ leads to a position
$S \cup \{n\}$ that the next player loses;
\item the player to move from $S$ \emph{loses} if every legal $n \ne 1$ leads to a position that the
next player wins.
\end{itemize}
Naming $1$ loses at once, so it is never a winning move. If $1$ is the only legal move, the player to
move loses. A first move $n$ is \emph{winning} if $n$ is legal after $\emptyset$, $n \ne 1$, and the
second player loses from $\{n\}$. A number is \emph{composite} if it is greater than $1$ and not prime.

\section{The proof}

\begin{lemma}[Determinacy]\label{lem:det}
If every legal move after $S$ is smaller than $B$, then the player to move from $S$ either wins or loses.
\end{lemma}
\begin{proof}
Use induction on the number of legal moves below $B$. Naming a legal $n$ removes $n$ from the legal set
and adds nothing to it. If some legal $n \ne 1$ leads to a lost position, the mover wins. Otherwise
every successor position is determined by induction and is not lost, so it is won, and the mover loses.
\end{proof}

\begin{lemma}[Closure invariance]\label{lem:cl}
If $\langle S\rangle = \langle T\rangle$, then $S$ and $T$ have the same game value.
\end{lemma}
\begin{proof}
Legality depends only on $\langle S\rangle$, and
$\langle S\cup\{n\}\rangle = \langle n\rangle + \langle S\rangle$. Now use induction on the game value.
\end{proof}

\begin{lemma}[Symmetry]\label{lem:sym}
Let $\gcd(p,q)=1$ with $p,q>1$, and let $t = pq-p-q$. If $x \notin \langle p,q\rangle$, then $x \le t$
and $t-x \in \langle p,q\rangle$.
\end{lemma}
\begin{proof}
The number $t$ is the Frobenius number of $\langle p,q\rangle$: the largest integer not in
$\langle p, q\rangle$. This is Mathlib's \texttt{frobeniusNumber\_pair}. Hence $x \le t$, and
$x+pq > t$ lies in $\langle p,q\rangle$. Write $x+pq = aq+bp$ with $a,b\ge 0$. If $a \ge p$, then
$x=(a-p)q+bp \in \langle p,q\rangle$, which is impossible. So $a \le p-1$, and likewise $b \le q-1$.
Then
\[
t-x = 2pq - p - q - (aq+bp) = (p-1-a)\,q + (q-1-b)\,p \in \langle p,q\rangle. \qedhere
\]
\end{proof}

\begin{theorem}[Hutchings]\label{thm:h}
If $p\ge 5$ is prime, then after the opening $p$ the second player loses.
\end{theorem}
\begin{proof}
Let the second player name a legal $q\ne 1$. Then $q>1$ and $p\nmid q$, so $\gcd(p,q)=1$. Let
$t=pq-p-q = (p-1)(q-1)-1 \ge 3$. By Lemma~\ref{lem:sym}, $t$ is legal after $\{p,q\}$, and every legal
move after $T=\{p,q,t\}$ is at most $t$. By Lemma~\ref{lem:det}, the player to move from $T$ either loses
or wins.
\begin{itemize}[nosep]
\item If the mover loses from $T$, the first player names $t$ and wins.
\item Otherwise there is a legal $x\ne1$ after $T$ such that the mover loses from $T\cup\{x\}$. The
number $x$ is also legal after $\{p,q\}$. By Lemma~\ref{lem:sym}, $t = x+(t-x) \in \langle p,q,x\rangle$,
so $\langle p,q,x\rangle = \langle p,q,t,x\rangle$. By Lemma~\ref{lem:cl}, the mover loses from
$\{p,q,x\}$, so the first player names $x$ and wins.
\end{itemize}
In both cases the first player wins after $q$, so the second player loses from $\{p\}$.
\end{proof}

\begin{corollary}
(C1) holds, but (C2) and (C3) fail.
\end{corollary}
\begin{proof}
The first move $5$ is legal, it is not $1$, and by Theorem~\ref{thm:h} it is winning. So the first
player wins, which is (C1). But $5$ is prime, hence not composite, so (C2) fails. The prime first move
$5$ is winning, and after it the second player has no winning response (a position is not both won and
lost), so both phrasings of (C3) fail.
\end{proof}

\paragraph{Literature.} Wikipedia, \emph{Sylver coinage},
\url{https://en.wikipedia.org/wiki/Sylver_coinage}: ``Hutchings's Theorem states that any of the prime
numbers 5, 7, 11, 13, ..., wins as a first move''. The strategy-stealing argument with the ``ender''
$t$ is also described there. The proof above is self-contained.

\paragraph{Readings.} The rules are the standard ones: legality is non-membership in the generated
submonoid, and naming $1$ loses. We formalize (C3) in both of its phrasings, ``no prime first move is
winning'' and ``after every prime opening the second player has a winning response''. Both are refuted
by $p=5$, and in fact by every prime $p\ge5$. Clause (C2) is refuted independently. The conjecture is
therefore false under any reading that keeps (C2) or (C3).

\section{Lean 4 formalization}

The file \texttt{lean/Conjecture1213/Basic.lean} (Lean \texttt{v4.33.1}, Mathlib \texttt{v4.33.1})
reuses \texttt{Legal}, \texttt{Outcome}, \texttt{anyNumber} and \texttt{not\_win\_and\_lose} from the
accepted solution of Conjecture 00000008869 by feiyuceng06-prog (GPL-3.0), with credit. Its main
declarations, all in namespace \texttt{Submission00000001213}, are:
\begin{itemize}[nosep]
\item \texttt{Clause1}, \texttt{Clause2}, \texttt{Clause3}, \texttt{Clause3'} and \texttt{Conjecture}:
the statement;
\item \texttt{determined}: Lemma~\ref{lem:det};
\item \texttt{outcome\_congr}: Lemma~\ref{lem:cl};
\item \texttt{gap\_symm}: Lemma~\ref{lem:sym}, via Mathlib's \texttt{frobeniusNumber\_pair} and
\texttt{AddSubmonoid.mem\_closure\_pair};
\item \texttt{two\_generator\_win} and \texttt{prime\_opening\_wins}: Theorem~\ref{thm:h};
\item \texttt{hutchings}: every prime $p\ge5$ is a winning first move;
\item \texttt{clause1\_true}, \texttt{not\_clause2}, \texttt{not\_clause3} and
\texttt{not\_clause3'}.
\end{itemize}
The main theorem is
\begin{center}
\texttt{theorem conjecture\_00000001213\_false :} $\neg$\,\texttt{Conjecture}
\end{center}
\texttt{lake build} succeeds. \texttt{\#print axioms} for \texttt{conjecture\_00000001213\_false},
\texttt{not\_clause3}, \texttt{not\_clause3'}, \texttt{hutchings} and \texttt{clause1\_true} reports only
\texttt{propext}, \texttt{Classical.choice} and \texttt{Quot.sound}. There is no \texttt{sorry} and no
\texttt{native\_decide}.

\end{document}
